\rmvspace
\subsection{Verschiedene Funktionen}
\shortremark[Ungleichungen] Zur Erinnerung: Richtung wechselt bei Muliplikation mit (Division durch) negative Zahl.

\subsubsection{Logarithmen}
Hier ist wieder $\log = \ln$ (wenn keine Basis gegeben)
\begin{itemize}
    \item \textit{(Basiswechsel)} $\log_a(x) = \frac{\ln(x)}{\ln(a)}$
    \item \textit{(Potenzen)} $\log_a(x^y) = y\log_a(x)$
    \item \textit{(Div, Mul)} $\log_a(x \cdot (\div) y) = \log_a(x) +(-) \log_a(y)$
    \item $\log_a(1) = 0 \smallhspace \forall a \in \N$
\end{itemize}


\subsubsection{Trigonometrie}
$\cot(\xi) = \displaystyle\frac{\cos(\xi)}{\sin(\xi)}, \tan(\xi) = \frac{\sin(\xi)}{\cos(\xi)}$

$\sinh(x) := \frac{e^x - e^{-x}}{2} : \R \rightarrow \R$,
$\cosh(x) := \frac{e^x + e^{-x}}{2} : \R \rightarrow [1, \infty]$,
$\cosh(x) := \frac{\sinh(x)}{\cosh(x)} = \frac{e^x - e^{-x}}{e^x + e^{-x}} : \R \rightarrow [-1, 1]$

\begin{enumerate}
    \item $\cos(x) = \cos(-x)$ \trand $\sin(-x) = -\sin(x)$
    \item $\cos(\pi - x) = -\cos(x)$ \trand $\sin(\pi - x) \sin(x)$
    \item $\sin(x + w) = \sin(x) \cos(w) + \cos(x) \sin(w)$
    \item $\cos(x + w) = \cos(x) \cos(w) - \sin(x) \sin(w)$
    \item $\cos(x)^2 + \sin(x)^2 = 1$
    \item $\sin(2x) = 2 \sin(x) \cos(x)$
    \item $\cos(2x) = \cos(x)^2 - \sin(x)^2$
\end{enumerate}


\shade{teal}{\tr{Values of trigonometric functions}{Werte der trigonometrischen Funktionen}}
\begin{tables}{ccccc}{° & rad              & $\sin(\xi)$          & $\cos(\xi)$           & $\tan(\xi)$}
              0°        & $0$              & $0$                  & $1$                   & $1$                   \\
              \hline
              30°       & $\frac{\pi}{6}$  & $\frac{1}{2}$        & $\frac{\sqrt{3}}{2}$  & $\frac{\sqrt{3}}{2}$  \\
              \hline
              45°       & $\frac{\pi}{4}$  & $\frac{\sqrt{2}}{2}$ & $\frac{\sqrt{2}}{2}$  & $1$                   \\
              \hline
              60°       & $\frac{\pi}{3}$  & $\frac{\sqrt{3}}{3}$ & $\frac{1}{2}$         & $\sqrt{3}$            \\
              \hline
              90°       & $\frac{\pi}{2}$  & $1$                  & $0$                   & $\varnothing$         \\
              \hline
              120°      & $\frac{2\pi}{3}$ & $\frac{\sqrt{3}}{2}$ & $-\frac{1}{2}$        & $-\sqrt{3}$           \\
              \hline
              135°      & $\frac{3\pi}{4}$ & $\frac{\sqrt{2}}{2}$ & $-\frac{\sqrt{2}}{2}$ & $-1$                  \\
              \hline
              150°      & $\frac{5\pi}{6}$ & $\frac{1}{2}$        & $-\frac{\sqrt{3}}{2}$ & $-\frac{\sqrt{3}}{2}$ \\
              \hline
              180°      & $\pi$            & $0$                  & $-1$                  & $0$                   \\
\end{tables}


\subsubsection{Reihen}
$\prod_{k = 1}^{n} \vartheta (1 - \vartheta)^{x_k - 1} = \vartheta^n (1 - \vartheta)^{S - n}$ mit $S = \sum_{k = 1}^{n} x_k$

$\sum_{k = 1}^n k = \frac{n \cdot (n + 1)}{2}$; $\sum_{k = 1}^{n} k^2 = \frac{n \cdot (n + 1) \cdot (2n + 1)}{6}$;

$\sum_{k = 1}^{n} k^3 = \left( \frac{n \cdot (n + 1)}{2} \right)^2$;
$\sum_{k = 1}^{n} \frac{1}{k^2} = \frac{\pi^2}{6}$

Arithmetic series: $s_n = \frac{n \cdot (2a_1 + (n - 1) \cdot d)}{2}$, mit $d$ Differenz

Geometric series: $s_n = a_1 \frac{1 - q^n}{1 - q}$, mit $q$ Quotient
